\chapter{ماذا تحسب عصبونات \nt{GLUT}}
\label{sec:glutcomputer}
\begin{chaptergoals}
\item كيف يصنع العصبون المكامل التسريبي ذو العتبة \foreignlanguage{english}{OR} وكاشف تزامن، أي \foreignlanguage{english}{AND} في الزمن؛
\item لماذا لا تحسب عصبونات \nt{GLUT} وحدها إلا دوالّ رتيبة، فلا تستطيع فعاليةٌ أن توقف فعالية؛
\item كيف تتيح أزواج أسلاك ”التشغيل/الإطفاء“ لعصبونات \nt{GLUT} حساب أي دالة منطقية وإعلان الجاهزية؛
\item كيف تحوّل خطوط التأخير فرق الزمن إلى موضع، وكيف تصنع التغذية الراجعة الذاكرة والمؤقتات.
\end{chaptergoals}

\section{قواعد اللعبة}

رأينا في الفصل~\ref{sec:neurontypes} أن الغالبية الساحقة من عصبونات الدماغ من نوع \nt{GLUT}: فهي تثير فقط، وأوزانها موجبة فقط. لنأخذ من هذه العصبونات ما نشاء ونسأل: ماذا تستطيع مثل هذه الشبكة أن تحسب، إذا لم نسمح لأنفسنا إلا بما يحدث في الكائن الحي؟

وليس في الكائن الحي مولّد نبضات ساعة يبدّل كل العناصر في آن واحد. كل عصبون يعمل على حدة، في الزمن المتصل: تأتي النبضات وتذهب متى اتفق، بتأخيرات مختلفة. ثمة \emph{مداخل}، هي العصبونات الحسية التي تنطلق بالضوء أو الصوت أو الضغط، وثمة \emph{مخارج}، هي العصبونات التي تتحكم في العضلات. وبينهما شبكة، لا يخبرها أحد متى تبدأ الحساب ومتى تنهيه. بالنسبة لمهندس الإلكترونيات هذه دارة \emph{لا تزامنية} عناصرها ذوات تأخيرات غير معروفة بدقة مسبقًا.

لنأخذ للعصبون أبسط النماذج التي تعمل بأمانة في الزمن المتصل. الجهد $V$ على غشاء العصبون يساوي صفرًا في الراحة. كل نبضة واردة من العصبون $j$ ترفعه قفزةً بمقدار $w_j\ge0$، ثم ينساب الجهد من تلقاء نفسه عائدًا إلى الصفر بثابت زمني $\tau$:
\begin{empheq}[box=\keybox]{equation}
\tau\,\frac{\dd V}{\dd t} \;=\; -V \;+\; \tau\sum_j w_j \sum_k \delta\bigl(t-t_j^{(k)}-d_j\bigr), \qquad w_j\ge 0 .
\label{eq:glutneuron}
\end{empheq}
هنا $t_j^{(k)}$ لحظات نبضات العصبون $j$، و$d_j$ التأخير على الطريق منه، و$\delta$ دالة دلتا، وهي التي تعني القفزة اللحظية. حين يبلغ $V$ العتبة $\theta$ يطلق العصبون نبضة، ويُصفَّر $V$، ولا يستطيع العصبون أن ينطلق ثانية مدة جزء من الألف من الثانية تقريبًا. الأرقام النموذجية: $\tau$ من أجزاء من جزء الألف من الثانية إلى عشرين جزءًا من الألف في العصبونات المختلفة، والتأخيرات من أجزاء من جزء الألف إلى عشرات الأجزاء من الألف من الثانية. هذا \emph{عصبون مكامل تسريبي}: دارة RC بعتبة. وهو تبسيط مقصود: ففي عصبونات القشرة تطلق فروع المداخل نفسها نبضات محلية~\cite{Smith2013}، ويتصرف العصبون الواحد أقرب إلى شبكة صغيرة متعددة الطبقات (الفصل~\ref{sec:synthesis}). لكن دارة RC واحدة تكفي لأسئلة هذا الفصل.

الفرق الرئيسي عن البوابة المنطقية هو التسريب: لا يتذكر العصبون النبضة إلى الأبد، بل نحو $\tau$، ولا يُجمع إلا ما وصل ضمن هذه النافذة.

\section{\foreignlanguage{english}{OR} و\foreignlanguage{english}{AND}}

لنبدأ بالبوابات. إذا رفعت نبضة واحدة الجهد فوق العتبة، $w\ge\theta$، انطلق العصبون مع كل نبضة من أي مدخل. هذه بوابة \emph{\foreignlanguage{english}{OR}}.

أما \foreignlanguage{english}{AND} فأكثر إثارة. ليكن $w<\theta<2w$: نبضة واحدة لا تكفي، ويلزم نبضتان. لكن السؤال ”هل وصلت النبضتان كلتاهما“ يصبح بلا معنى ما لم نحدد \emph{ضمن أي زمن}. لنفترض أن النبضة الأولى وصلت في اللحظة $0$، والثانية بعد زمن $\Delta$. عند وصول الثانية يكون قد بقي من الأولى $we^{-\Delta/\tau}$، وينطلق العصبون إذا كان $we^{-\Delta/\tau}+w\ge\theta$. ومن هنا نافذة التزامن:
\begin{empheq}[box=\keybox]{equation}
\Delta \;\le\; \Delta_{\max} \;=\; \tau\,\ln\frac{w}{\theta-w} .
\label{eq:window}
\end{empheq}
هذه ليست بوابة \foreignlanguage{english}{AND} منطقية، بل \emph{كاشف تزامن}: \foreignlanguage{english}{AND} في الزمن (الشكل~\ref{fig:coincidence}). عند $\theta=1.5w$ تساوي النافذة $\tau\ln2\approx0.7\tau$. وفي عصبون قشري بـ$\tau=10\ms$ هذا نحو سبعة أجزاء من الألف من الثانية. أما في عصبونات الجهاز السمعي، حيث $\tau$ أقل من جزء من الألف من الثانية، فتضيق النافذة إلى أجزاء من جزء الألف.

\begin{figure}[t]
\centering
\begin{tikzpicture}[>=Stealth,x=0.9cm,y=1.6cm]
\foreach \dx/\lab/\c [count=\i] in {0.5/{تزامنتا}/red!70!black, 2.6/{لم تتزامنا}/blue!60!black}{
 \begin{scope}[xshift={(\i-1)*7.2cm}]
  \draw[->,gray!70!black] (0,0) -- (6.3,0) node[below left,black] {\small $t$};
  \draw[->,gray!70!black] (0,0) -- (0,1.75) node[left,black] {\small $V$};
  \draw[densely dashed,gray!60!black] (0,1.5) -- (6.0,1.5) node[right,black] {\small $\theta$};
  \draw[very thick,\c] (0,0) -- (1,0) -- (1,1)
      plot[domain=1:1+\dx,samples=30] (\x,{exp(-(\x-1)/1.5)});
  \pgfmathsetmacro{\v}{exp(-\dx/1.5)+1}
  \draw[very thick,\c] (1+\dx,{exp(-\dx/1.5)}) -- (1+\dx,{min(\v,1.5)});
  \ifdim\v pt<1.5pt
    \draw[very thick,\c] plot[domain=1+\dx:6,samples=40] (\x,{\v*exp(-(\x-1-\dx)/1.5)});
  \else
    \draw[very thick,\c] (1+\dx,1.5) -- (1+\dx,0) -- (6,0);
    \draw[->,thick,\c] (1+\dx,1.55) -- (1+\dx,1.72) node[above,font=\scriptsize] {نبضة};
  \fi
  \draw[->,thick] (1,-0.35) -- (1,-0.08); \draw[->,thick] (1+\dx,-0.35) -- (1+\dx,-0.08);
  \draw[decorate,decoration={brace,amplitude=3pt,mirror}] (1,-0.42) -- (1+\dx,-0.42) node[midway,below=3pt] {\scriptsize $\Delta$};
  \node[\c] at (3.5,-0.95) {\small \lab};
 \end{scope}}
\end{tikzpicture}
\caption{كاشف التزامن، والعتبة $\theta=1.5w$. على اليسار وصلت النبضة الثانية بعد $\Delta<\Delta_{\max}$، فبلغ المجموع العتبة وانطلق العصبون. وعلى اليمين $\Delta>\Delta_{\max}$: لم يبق من النبضة الأولى شيء يُذكر، فصمت العصبون.}
\label{fig:coincidence}
\end{figure}

وثمة طريقة أخرى للحصول على \foreignlanguage{english}{AND}، عبر المدة لا عبر التوقيت الدقيق. ما دام الضوء مضاءً يطلق العصبون الحسي نبضات متكررة ومنتظمة؛ فإذا لم يكفِ العتبةَ مدخلٌ واحد كهذا وكفاها مدخلان، عمل العصبون ما دام المدخلان كلاهما نشطين. هكذا تحسب الشبكة \emph{بالمستويات}، ولا يهم التوقيت الدقيق لوصول النبضات. ومعظم ما يلي سيكون على هذا النحو.

\section{ما لا يمكن فعله}

لنحاول الآن صنع بوابة \foreignlanguage{english}{NOT}: عصبون يعمل حين يصمت المدخل. لن ننجح: من دون نبضات دخل لا توجد في~\eqref{eq:glutneuron} قفزة واحدة، ويبقى الجهد عند الصفر، تحت العتبة. وماذا عن شبكة من عصبونات كثيرة؟ لا أيضًا، ويُبرهَن ذلك دفعة واحدة لكل الشبكات.

أيسر ما يكون الحديث بالترددات. لتكن $r_i$ تردد نبضات العصبون $i$، و$I_i$ الوارد إليه من المداخل، و$f_i$ منحنى نقله: كيف يتعلق التردد بمجموع الوارد. هذا المنحنى في العصبون المكامل غير متناقص: وارد أكبر يعني نبضات أكثر تواترًا. والترددات في شبكة بلغت الاتزان تحقق المعادلات
\begin{equation}
r_i \;=\; f_i\Bigl(\sum_j w_{ij}\,r_j + I_i\Bigr), \qquad w_{ij}\ge 0 .
\label{eq:ratenet}
\end{equation}

\begin{keyblock}
\begin{proposition}[الرتابة]
\label{prop:monotone}
لتكن الشبكة~\eqref{eq:ratenet} مؤلفة من عصبونات \nt{GLUT} فقط، ولتبدأ من صمت تام. إذا قُوّيت المداخل فلن ينقص التردد المستقر لأي عصبون. كل ما تحسبه هذه الشبكة دوال \emph{رتيبة} (غير متناقصة) في المداخل.
\end{proposition}
\end{keyblock}

البرهان. لنبدأ من الصمت، $r^{(0)}=0$، ولنعوّض الترددات في الطرف الأيمن من~\eqref{eq:ratenet}: $r^{(k+1)}=F\bigl(r^{(k)}\bigr)$. الطرف الأيمن $F$ غير متناقص في أي من متغيراته، لأن الأوزان غير سالبة والدوال $f_i$ غير متناقصة. وبما أن $r^{(1)}\ge0=r^{(0)}$، فبالاستقراء $r^{(k+1)}\ge r^{(k)}$: الترددات تزداد فقط وتصل إلى أصغر حل لـ~\eqref{eq:ratenet}. وإذا استُبدلت بالمداخل $I$ مداخل أكبر $I'$، كان $r^{(k)}(I)\le r^{(k)}(I')$ في كل خطوة، وفي النهاية كذلك. الشبكة الحقيقية تبلغ الاتزان لا بخطوات بل باستمرار، لكن الأمر نفسه صحيح لها: مع الوصلات الموجبة تبقى المتباينة بين تشغيلين، إن تحققت في البداية، متحققة طوال الوقت.

أما للنبضات المفردة فالقضية لا تصح حرفيًا: قد تجعل نبضة دخل زائدة العصبونَ ينطلق أبكر قليلًا، فتضيع نبضته التالية بسبب جزء الألف من الثانية الذي لا يستجيب فيه. لكن الأثر ضعيف وغير موثوق، ولا يصلح أساسًا للحساب. أما للترددات فالرتابة صارمة.

والنتائج هي نتائج أي دارة بلا عواكس. لا \foreignlanguage{english}{NOT}. ولا \foreignlanguage{english}{XOR}: إضافة مدخل ثانٍ لا تستطيع إطفاء المخرج. والأسوأ: \emph{لا يمكن إيقاف الفعالية بفعالية}، بل يمكن فقط إزالة ما يغذيها.

\section{سلكان: التشغيل والإطفاء}

المخرج من هذا المأزق يوحي به الكائن الحي نفسه. في كثير من أعضاء الحس ينتقل المنبه عبر مجموعتين من العصبونات: ففي البصر تستجيب بعض الخلايا لاشتداد الضوء، وأخرى لخفوته~\cite{Schiller1992}. لنسمها عصبونات \emph{التشغيل} وعصبونات \emph{الإطفاء}. فبدل السلك الواحد $x$ تتلقى الشبكة الزوج $(x,\bar x)$، والغياب الذي لا تستطيع حسابه يبلغها به المستشعر نفسه.

مع زوج من الأسلاك تصبح \foreignlanguage{english}{NOT} مجانية: يكفي أن نبدّل السلكين. و\foreignlanguage{english}{AND} و\foreignlanguage{english}{OR} تُصنع كل منهما بعصبونين:
\begin{empheq}[box=\keybox]{equation}
\begin{aligned}
x\wedge y \;&\to\; \bigl(\;x\wedge y,\;\; \bar x\vee\bar y\;\bigr),\\
x\vee y \;&\to\; \bigl(\;x\vee y,\;\; \bar x\wedge\bar y\;\bigr).
\end{aligned}
\label{eq:dualrail}
\end{empheq}
كل العمليات في الطرف الأيمن رتيبة، فلا تُنتهك القضية~\ref{prop:monotone}.

وللشيفرة ثنائية السلك ميزة ثانية للدارة اللاتزامنية، أهم حتى من الأولى. للزوج ثلاث حالات: ”نعم“، و”لا“، و\emph{”لا جواب بعد“} حين يصمت السلكان. يعرف المستقبِل أن الحساب انتهى لا من ساعة، بل من الإشارة نفسها: في كل زوج اشتغل سلك واحد بالضبط. وبين منبه وآخر تصمت المستشعرات، فتعود الشبكة إلى الصمت، مستعدة للمرة التالية. وعلى هذه الحيلة تُبنى في الإلكترونيات اللاتزامنية دارات تعمل صحيحةً مهما كانت تأخيرات عناصرها~\cite{Sparso2001}.

\begin{keyblock}
\begin{proposition}[الحساب اللاتزامني]
\label{prop:dualrail}
إذا ورد كل مدخل على زوج من الأسلاك ”تشغيل/إطفاء“، استطاعت شبكة من عصبونات \nt{GLUT} أن تحسب أي دالة منطقية في المداخل. ويأتي الجواب أيضًا على أزواج من الأسلاك، ويُعرف جاهزيته من أن سلكًا واحدًا اشتغل في كل زوج. لا حاجة إلى ساعة لذلك. والثمن: عدد عصبونات أكبر بمرتين تقريبًا.
\end{proposition}
\end{keyblock}

مثال ذلك \foreignlanguage{english}{XOR}: ”تغيّر منبه واحد بالضبط من منبهين“. السلك المباشر للجواب هو $(x\wedge\bar y)\vee(\bar x\wedge y)$، والسلك المعاكس $(x\wedge y)\vee(\bar x\wedge\bar y)$. هذه ستة عصبونات، لا يحتاج أي منها إلى التثبيط (الشكل~\ref{fig:dualxor}).

\begin{figure}[t]
\centering
\begin{tikzpicture}[>=Stealth,x=1cm,y=1cm,
  nrn/.style={circle,draw,very thick,fill=blue!6,minimum size=0.75cm,inner sep=0pt,font=\small},
  inp/.style={font=\small}]
\node[inp] (X)  at (0,3.3) {$x$};
\node[inp] (Xb) at (0,2.2) {$\bar x$};
\node[inp] (Y)  at (0,1.1) {$y$};
\node[inp] (Yb) at (0,0)   {$\bar y$};
\node[nrn] (A1) at (3.2,3.3) {$2$};
\node[nrn] (A2) at (3.2,2.2) {$2$};
\node[nrn] (A3) at (3.2,1.1) {$2$};
\node[nrn] (A4) at (3.2,0)   {$2$};
\node[nrn] (O)  at (7.2,2.75) {$1$};
\node[nrn] (Ob) at (7.2,0.55) {$1$};
\foreach \s/\t in {X/A1,Yb/A1,Xb/A2,Y/A2,X/A3,Y/A3,Xb/A4,Yb/A4}{\draw[->,thick,blue!60!black] (\s) -- (\t);}
\draw[->,thick,blue!60!black] (A1) -- node[pos=0.45,above,sloped,font=\scriptsize,black] {$x\wedge\bar y$} (O);
\draw[->,thick,blue!60!black] (A2) -- node[pos=0.45,below,sloped,font=\scriptsize,black] {$\bar x\wedge y$} (O);
\draw[->,thick,blue!60!black] (A3) -- node[pos=0.45,above,sloped,font=\scriptsize,black] {$x\wedge y$} (Ob);
\draw[->,thick,blue!60!black] (A4) -- node[pos=0.45,below,sloped,font=\scriptsize,black] {$\bar x\wedge\bar y$} (Ob);
\draw[->,very thick,red!70!black] (O) -- (8.6,2.75) node[right,black] {$x\oplus y$: ”نعم“};
\draw[->,very thick,red!70!black] (Ob) -- (8.6,0.55) node[right,black] {$\overline{x\oplus y}$: ”لا“};
\node[below,font=\small,gray!40!black] at (3.2,-0.55) {\foreignlanguage{english}{AND}: العتبة $2$};
\node[below,font=\small,gray!40!black] at (7.2,-0.55) {\foreignlanguage{english}{OR}: العتبة $1$};
\node[below,font=\small,gray!40!black] at (0,-0.55) {أزواج الأسلاك};
\end{tikzpicture}
\caption{بوابة \foreignlanguage{english}{XOR} بالشيفرة ثنائية السلك، من عصبونات \nt{GLUT} وحدها. العدد في الدائرة هو العتبة بوحدات وزن مدخل واحد. أربعة عصبونات \foreignlanguage{english}{AND} تتعرف على التوليفات الأربع للمداخل، وعصبونا \foreignlanguage{english}{OR} يجمعانها في زوج أسلاك الجواب.}
\label{fig:dualxor}
\end{figure}

\section{التأخيرات بدل الساعة}

للحساب بالزمن تكفي التأخيرات في الأسلاك وكواشف التزامن. وأفضل مثال طريقة الدماغ في تحديد الجهة التي أتى منها الصوت.

يصل الصوت الآتي من مصدر جانبي إلى الأذن القريبة قبل البعيدة. ويتوقف الفرق على الاتجاه: من الصفر حين يكون المصدر أمامنا مباشرة، إلى $0.22\,\text{m}/343\,\text{m/s}\approx0.6\ms$ عند الإنسان حين يكون على الجانب تمامًا. يميّز الدماغ فرقًا من نحو عشرة \emph{أجزاء من المليون} من الثانية~\cite{KlumppEady1956,Grothe2010}، مع أن العصبونات نفسها لا تنطلق بدقة أفضل من أجزاء من جزء الألف من الثانية. كيف؟

لنمدّ من الأذن اليسرى سلكًا على طول صف من كواشف التزامن من اليسار إلى اليمين، ومن الأذن اليمنى سلكًا من اليمين إلى اليسار (الشكل~\ref{fig:itd}). تجري الإشارة في السلك بسرعة $v$. إلى الكاشف الواقع على بعد $x$ من الطرف الأيسر لصف طوله $L$، تستغرق الإشارة اليسرى زمنًا $x/v$، واليمنى $(L-x)/v$. فإذا وصل الصوت إلى الأذن اليمنى متأخرًا بـ$\delta t$، التقت الإشارتان في آن واحد عند النقطة التي
\begin{empheq}[box=\keybox]{equation}
\frac{x}{v} \;=\; \frac{L-x}{v} + \delta t
\quad\Longrightarrow\quad
x \;=\; \frac{L}{2} + \frac{v\,\delta t}{2} .
\label{eq:itd}
\end{empheq}
لن ينطلق إلا الكاشف الذي بلغته الإشارتان ضمن النافذة~\eqref{eq:window}. ورقم الكاشف المنطلق هو اتجاه المصدر. لقد تحوّل الفرق في الزمن إلى \emph{موضع}.

\begin{figure}[t]
\centering
\begin{tikzpicture}[>=Stealth,
  nrn/.style={circle,draw,very thick,fill=blue!6,minimum size=0.7cm,inner sep=0pt}]
\foreach \k in {0,...,6}{\node[nrn] (D\k) at (1.4*\k,0) {\scriptsize $2$};}
\node[nrn,fill=red!15] at (D4) {\scriptsize $2$};
\draw[->,thick,blue!60!black] (-1.4,0.9) node[left,black] {\small الأذن اليسرى} -- (9.2,0.9);
\draw[->,thick,orange!85!black] (9.8,-0.9) node[right,black] {\small الأذن اليمنى} -- (-0.8,-0.9);
\foreach \k in {0,...,6}{
  \draw[->,blue!60!black] (1.4*\k,0.9) -- (D\k.north);
  \draw[->,orange!85!black] (1.4*\k,-0.9) -- (D\k.south);}
\draw[->,thick,red!70!black] (D4.north east) ++(0.1,0.05) -- ++(0.5,0.45) node[above right,font=\small,black] {انطلق};
\node[font=\small,gray!40!black] at (4.2,-1.5) {يزداد التأخير $\longleftarrow$ \hspace{2.2cm} $\longrightarrow$ يزداد التأخير};
\pic[scale=1.4] at (-2.3,1.75) {speaker};
\node[font=\scriptsize,align=center] at (-2.3,2.35) {المصدر على اليسار};
\end{tikzpicture}
\caption{تحديد اتجاه الصوت. تجري إشارتا الأذنين إحداهما نحو الأخرى؛ وكل دائرة كاشف تزامن بعتبة $2$. إذا وصل الصوت إلى الأذن اليمنى متأخرًا، التقت الإشارتان يمين المنتصف، بحسب الصيغة~\eqref{eq:itd}. مخرج الشبكة هو رقم العصبون المنطلق. والمصدر هنا على اليسار: وصل الصوت إلى الأذن اليسرى أولًا.}
\label{fig:itd}
\end{figure}

لا ساعة هنا، ولا تثبيط، ولا حتى شيفرة ثنائية السلك: لا تجيب الشبكة بـ”نعم“ أو ”لا“، بل تشير إلى عصبون واحد من بين كثيرين. ويبدو أن دارة كهذه من خطوط التأخير وكواشف التزامن تعمل فعلًا في جذع الدماغ السمعي لدى الطيور~\cite{Carr1990}. ودقة أجزاء المليون من الثانية لا تأتي من دقة كل عصبون، بل من كثرة الكواشف ونظر كل منها إلى تأخيره الخاص. ولم يُعثر لدى الثدييات على صف كهذا من التأخيرات: هناك يبدو أن المهمة نفسها تحلها كواشف تزامن يضبطها تثبيط محسوب التوقيت بدقة من عصبونات \nt{GLYC}~\cite{Brand2002,Grothe2010}. وكيف يضيّق التثبيط نافذة التزامن سنبحثه في الفصل~\ref{sec:gaba} على مثال \nt{GABA}؛ والغليسين يثبط بالطريقة نفسها، إذ يفتح لدى المستقبِل قنوات الكلور.

\section{الذاكرة}

يحتاج الحاسوب إلى ذاكرة. والذاكرة في شبكة كهذه فعالية تغذي نفسها بنفسها. لنأخذ مجموعة من عصبونات \nt{GLUT} متصلة بعضها ببعض اتصالًا قويًا، ولنصفها بتردد واحد $r$ (الشكل~\ref{fig:recurrentgroup}أ). في الاتزان $r=f(wr+I)$، حيث $w$ قوة التغذية الراجعة للمجموعة على نفسها، و$I$ الوارد من الخارج. لنحل هذه المعادلة بيانيًا بتقاطع المنحنى $f(wr+I)$ مع المستقيم $r$ (الشكل~\ref{fig:bistable}).

\begin{figure}[t]
\centering
% (b): tau dr/dt = -r + f(w r + I), f as in fig:bistable, w=1, I=0.6; Euler dt=0.001 tau.
\begin{tikzpicture}[>=Stealth,
  nrn/.style={circle,draw,thick,fill=blue!6,minimum size=0.5cm,inner sep=0pt}]
\begin{scope}
\node[font=\small] at (-0.5,1.55) {(أ)};
\foreach \k in {1,...,6}{\node[nrn] (n\k) at ({1.4+1.0*cos(60*\k)},{1.0*sin(60*\k)}) {};}
\foreach \a/\b in {1/2,2/3,3/4,4/5,5/6,6/1,1/4,2/5,3/6}{\draw[<->,blue!60!black] (n\a) -- (n\b);}
\foreach \k in {2,3,4}{\draw[->,thick,green!45!black] ($(n\k)+(-0.95,0)$) -- (n\k);}
\node[font=\small,green!45!black] at (-0.35,-0.45) {$I$};
\node[font=\Large] at (3.1,0) {$\approx$};
\node[nrn,minimum size=0.85cm,font=\small] (R) at (4.35,-0.1) {$r$};
\draw[->,thick,blue!60!black] (R) edge[loop above,looseness=6] node[above,font=\small] {$w$} (R);
\draw[->,thick,green!45!black] (3.55,-0.95) node[left,font=\small] {$I$} -- (R);
\end{scope}
\begin{scope}[xshift=6.3cm,y=2.6cm,x=1.05cm]
\node[font=\small] at (-0.6,1.25) {(ب)};
\draw[->,gray!70!black] (0,0) -- (7.3,0) node[below left,font=\small,black] {$t/\tau$};
\draw[->,gray!70!black] (0,0) -- (0,1.12) node[left,font=\small,black] {$r$};
\foreach \t in {0,2,4,6}{\draw[gray!60!black] (\t,-0.02) -- (\t,0.02); \node[below,font=\scriptsize] at (\t,0) {\t};}
\draw[densely dotted,gray!60!black] (0,0.5) -- (7,0.5) node[right,font=\scriptsize,black,align=left] {عتبة\\الكتابة};
\fill[green!45!black,opacity=0.25] (1,-0.2) rectangle (2,-0.13);
\fill[gray!60!black,opacity=0.35] (1,-0.29) rectangle (1.5,-0.22);
\draw[very thick,blue!60!black] plot coordinates {(0.0,0.002) (0.1,0.002) (0.2,0.002) (0.3,0.002) (0.4,0.002) (0.5,0.002) (0.6,0.002) (0.7,0.002) (0.8,0.002) (0.9,0.002) (1.0,0.002) (1.1,0.083) (1.2,0.165) (1.3,0.242) (1.4,0.313) (1.5,0.378) (1.6,0.437) (1.7,0.491) (1.8,0.539) (1.9,0.583) (2.0,0.623) (2.1,0.643) (2.2,0.665) (2.3,0.688) (2.4,0.71) (2.5,0.732) (2.6,0.753) (2.7,0.773) (2.8,0.792) (2.9,0.81) (3.0,0.826) (3.2,0.855) (3.4,0.88) (3.6,0.9) (3.8,0.917) (4.0,0.931) (4.4,0.953) (4.8,0.967) (5.2,0.977) (5.6,0.984) (6.0,0.988) (6.5,0.992) (7.0,0.994)};
\draw[very thick,gray!60!black,densely dashed] plot coordinates {(0.0,0.002) (0.1,0.002) (0.2,0.002) (0.3,0.002) (0.4,0.002) (0.5,0.002) (0.6,0.002) (0.7,0.002) (0.8,0.002) (0.9,0.002) (1.0,0.002) (1.1,0.083) (1.2,0.165) (1.3,0.242) (1.4,0.313) (1.5,0.378) (1.6,0.358) (1.7,0.336) (1.8,0.313) (1.9,0.291) (2.0,0.269) (2.2,0.228) (2.4,0.191) (2.6,0.159) (2.8,0.133) (3.0,0.11) (3.2,0.091) (3.4,0.076) (3.6,0.063) (3.8,0.052) (4.0,0.043) (4.4,0.03) (4.8,0.021) (5.2,0.015) (5.6,0.011) (6.0,0.008) (6.5,0.006) (7.0,0.004)};
\node[font=\scriptsize,blue!60!black,anchor=west] at (3.3,0.8) {وارد لمدة $1\tau$: كُتب};
\node[font=\scriptsize,gray!50!black,anchor=west] at (2.3,0.3) {وارد لمدة $0.5\tau$: انطفأ};
\end{scope}
\end{tikzpicture}
\caption{مجموعة متصلة بنفسها اتصالًا قويًا. (أ)~عصبونات \nt{GLUT} كثيرة يثير بعضها بعضًا تتصرف كعصبون واحد بتردد $r$ يثير نفسه بقوة $w$؛ ويأتيه من الخارج وارد $I$. (ب)~تردد المجموعة مع الزمن عند $w=1$ وبالمنحنى $f$ نفسه الذي في الشكل~\ref{fig:bistable}. يُقدَّم الوارد $I=0.6$ طوال الزمن المعلَّم بالشريط تحت المحور. إذا دام $1\tau$، لحقت المجموعة أن تتجاوز النقطة غير المستقرة $r=0.5$، فتتقد وتبقى متقدة بعد انتهاء الوارد. وإذا دام $0.5\tau$ فقط، لم تبلغ تلك النقطة فتنطفئ عائدة: لكتابة بت يجب أن يدوم الوارد أكثر من $0.7\tau$ تقريبًا.}
\label{fig:recurrentgroup}
\end{figure}

\begin{figure}[t]
\centering
\begin{tikzpicture}[>=Stealth,x=5cm,y=3.2cm]
\draw[->,gray!70!black] (0,0) -- (1.1,0) node[below left,black] {\small $r$};
\draw[->,gray!70!black] (0,0) -- (0,1.1);
\draw[thick,gray!60!black] (0,0) -- (1.05,1.05) node[right,black] {\small $r$};
\draw[very thick,blue!60!black] plot[domain=0:1.05,samples=80] (\x,{1/(1+exp(-(\x-0.5)/0.08))});
\draw[very thick,red!70!black,densely dashed] plot[domain=0:1.05,samples=80] (\x,{1/(1+exp(-(\x+0.3-0.5)/0.08))});
\fill[blue!60!black] (0.002,0.002) circle (0.018);
\draw[blue!60!black,fill=white,thick] (0.5,0.5) circle (0.018);
\fill[blue!60!black] (0.998,0.998) circle (0.018);
\node[below right,font=\small] at (0.02,0.0) {مطفأ};
\node[right,font=\small] at (0.52,0.46) {غير مستقر};
\node[below,font=\small] at (0.99,0.96) {مشغَّل};
\node[anchor=east,font=\small,red!70!black] at (0.19,0.62) {$f(wr+I)$};
\node[anchor=west,font=\small,blue!60!black] at (0.45,0.1) {$f(wr)$};
\end{tikzpicture}
\caption{ذاكرة من مجموعة عصبونات \nt{GLUT} ذات تغذية راجعة قوية. المنحنى المتصل من دون وارد خارجي: تقاطعان مستقران مع المستقيم $r$، ”مطفأ“ و”مشغَّل“، وتقاطع غير مستقر بينهما. والوارد القصير $I$ (المنحنى المتقطع) لا يُبقي إلا التقاطع العلوي.}
\label{fig:bistable}
\end{figure}

إذا كانت التغذية الراجعة قوية كانت التقاطعات ثلاثة: السفلي هو الصمت، والعلوي هو المجموعة متقدة بكامل طاقتها، والأوسط غير مستقر. حصلنا على عنصر بحالتين مستقرتين: قلّاب (\foreignlanguage{english}{flip-flop}). وكتابة الواحد فيه سهلة: وارد قصير يرفع المنحنى، فيختفي التقاطع السفلي، وتتقد المجموعة، وتبقى متقدة بعد انتهاء الوارد (الشكل~\ref{fig:recurrentgroup}ب). أما الوارد القصير جدًا فلا يلحق أن يبلغ بالمجموعة النقطة غير المستقرة، فتنطفئ عائدة. هذه الفعالية الذاتية الاستمرار، التي تدوم ثوانيَ بعد زوال المنبه، تُرصد فعلًا في الدماغ، وتُعدّ أساس الذاكرة قصيرة الأمد~\cite{Wang2001}. لكنها ليست الطريقة الوحيدة لتذكر الثواني. يمكن أيضًا حفظ المكتوب في متغير بطيء في المشابك نفسها: بعد رشقة من النبضات يدوم التيسير، كما في المشبك ”الشَّرطة“ في الفصل~\ref{sec:code}، نحو ثانية، وقد تكاد الشبكة تصمت بين ومضات قصيرة، فهي ليست قلّابًا بل مكثفًا مشحونًا~\cite{Mongillo2008}. وهذه الفترات ”الصامتة“ أثناء الاحتفاظ بالذاكرة مرصودة~\cite{Stokes2015}، والتخزين من دون نبضات دائمة أرخص طاقةً. أما أي الطريقتين تستعمل القشرة وفي أي حالة، فموضع نقاش.

وكيف نمحو؟ بحسب القضية~\ref{prop:monotone} لا سبيل إلى ذلك بالفعالية؛ لا يبقى إلا أن \emph{نزيل} شيئًا. الطريقة الأولى إزالة الدعم: إذا كانت المجموعة لا تتقد إلا مع وارد من مجموعة أخرى، انطفأت الذاكرة معه. والثانية التعب. ففي المشابك الحقيقية ينضب مخزون الناقل العصبي عند العمل الطويل~\cite{Tsodyks1997}، وتنمو في العصبونات تيارات بطيئة تخفض الاستثارة. فتضعف التغذية الراجعة، ويختفي التقاطع العلوي، وتنطفئ المجموعة وحدها بعد ثوانٍ. والنتيجة ذاكرة تُشغَّل بإشارة ولا تُطفأ إلا بالزمن.

\section{المتتاليات}

بدل الساعة يمكن أن يكون لدينا \emph{مؤقت يطلقه حدث}. لنصل مجموعات من عصبونات \nt{GLUT} في سلسلة: الأولى تثير الثانية، والثانية الثالثة، وهكذا. تُشعل إشارة خارجية المجموعة الأولى، فتجري موجة الفعالية على طول السلسلة، كل حلقة بعد سابقتها بتأخير أو تأخيرين، ومرة واحدة فقط: فالعصبونات لا تستجيب مباشرة بعد النبضة، والموجة قد مضت.

تقيس هذه السلسلة الزمن منذ الحدث: إذا انطلقت الحلقة رقم $k$ فقد مضى منذ الإطلاق $k$ تأخيرات. ويمكن توصيل مخارجها بالعضلات للحصول على متتالية حركات تتكشف من تلقاء نفسها. لدى الطيور المغردة، أثناء الغناء، تنطلق عصبونات \nt{GLUT} مفردة في إحدى مناطق الدماغ بالتتابع الصارم، كل منها في لحظته من الأغنية، وهذا يشبه كثيرًا سلسلة كهذه~\cite{Hahnloser2002}.

وعلى خلاف الساعة، تصمت السلسلة ما لم تُطلق، وتعمل مرة واحدة، ولا تقيس الزمن إلا لمن وُصل بها. ولا يزال لا إيقاع مشتركًا للشبكة.

\section{ما الذي ينقص}

من عصبونات \nt{GLUT} وحدها، بلا تثبيط، يمكن صنع الكثير. لكن ثلاثة أشياء تنقص، وكلها بسبب القضية~\ref{prop:monotone}.

\emph{الاختيار.} يجب على الشبكة أن تختار اتجاهًا واحدًا، وفعلًا واحدًا. فإذا تساوى منبهان تقريبًا، انطلق في شبكة الشكل~\ref{fig:itd} المخرجان كلاهما، ولا شيء يكبت الأضعف لصالح الأقوى.

\emph{التصفير.} لا يمكن محو الذاكرة بإشارة.

\emph{الاستقرار.} في شبكة لا تنشأ فيها الوصلات بحسب مخطط مهندس، لا مفر من حلقات التغذية الراجعة الموجبة، وقد تتصاعد فيها الفعالية تصاعدًا متسلسلًا (الفصل~\ref{sec:neurontypes}). والكابح الوحيد هو التعب نفسه، البطيء والخشن.

والعلل الثلاث يعالجها دواء واحد: عصبون يطفئ النبضة بنبضة. هذا هو \nt{GABA}، والحاجة إليه ليست للحساب، بل للاختيار، والتصفير، وإبقاء الحساب تحت السيطرة.

\section{تمارين}

\begin{exercise}[\lvlA]
\label{ex:02:1}
\emph{صف الكواشف للصوت.} تجري إشارتا الأذنين في أسلاك الشكل~\ref{fig:itd} بسرعة $5\,\text{m/s}$، وأكبر فرق في الوصول $0.64\ms$. الكواشف عصبونات~\eqref{eq:glutneuron} بـ$\tau=0.5\ms$ و$\theta=1.5w$، مرتبة بحيث يميّز المتجاوران $10\,\mu\text{s}$. كم كاشفًا في الصف، وكم منها ينطلق لصوت واحد؟
\end{exercise}

\begin{exercise}[\lvlB]
\label{ex:02:2}
\emph{المداخل غير المتساوية تزيح النافذة.} كاشف تزامن~\eqref{eq:glutneuron} بـ$\tau=0.5\ms$ و$\theta=1.5$ يتلقى نبضة من كل من مدخلين بوزنين $1$ و$0.8$. أوجد نافذة التزامن ومركزها. بكم جزءًا من المليون من الثانية يخطئ صف الشكل~\ref{fig:itd} إذا كان مدخل إحدى الأذنين أضعف من الأخرى بـ$20\,\%$؟
\end{exercise}

\begin{exercise}[\lvlB]
\label{ex:02:3}
\emph{أي الشبكات لا تتذبذب.} الشبكة $\tau\dot r_i=-r_i+f_i\bigl(\textstyle\sum_jw_{ij}r_j+I_i\bigr)$ ذات الدوال $f_i$ غير المتناقصة و$w_{ij}\ge0$ عند $i\ne j$ رتيبة ولا تتذبذب تذبذبًا مستقرًا. برهن: بالتعويض $r_i\to\pm r_i$ تؤول إليها الشبكة التي في كل دورة فيها عدد زوجي من الوصلات المثبطة. هل يمكن أن يتذبذب التثبيط المتبادل بين مجموعتين؟ وحلقة \foreignlanguage{english}{\nt{GLUT}--\nt{GABA}}؟
\end{exercise}

\begin{exercise}[\lvlB]
\label{ex:02:4}
\emph{تصميم: جامع ثنائي السلك.} كل بت ينتقل على زوج من الأسلاك $(x,\bar x)$. ابنِ جامعًا كاملًا لثلاثة بتات يُخرج الزوجين $(S,\bar S)$ للمجموع و$(C,\bar C)$ للمنقول، من عصبونات \nt{GLUT}، مع بيان الأوزان والعتبات. اكتفِ بثمانية عصبونات. كم عصبونًا تتطلب دارة من عناصر \foreignlanguage{english}{AND} و\foreignlanguage{english}{OR} كما في الشكل~\ref{fig:dualxor}؟
\end{exercise}

\begin{exercise}[\lvlB]
\label{ex:02:5}
\emph{نمذجة: كم تدوم الكتابة.} مجموعة الشكل~\ref{fig:recurrentgroup}: $\tau\dot r=-r+f(r+I)$، $f(u)=1/\bigl(1+e^{-(u-0.5)/0.08}\bigr)$، وهي قبل الكتابة في الحالة السفلى؛ ويُقدَّم الوارد $I$ مدة $T$. أوجد عدديًا أصغر $T$ يكتب البت عند $I=0.6$، والعتبة $I_c$ التي لا يُكتب البت دونها مهما كان $T$.
\end{exercise}

\begin{exercise}[\lvlA]
\label{ex:02:6}
\emph{السلسلة مؤقتًا.} حلقة السلسلة $100$ عصبون من \nt{GLUT}، وتأخيرها $2\ms$ بتشتت مستقل $0.2\ms$. (أ)~كم عصبونًا يلزم لقياس $1\,\text{s}$، وما الخطأ النسبي؟ وكيف يتعلق بالفترة؟ (ب)~خطأ المؤقت الحقيقي $5\,\%$ لأي فترة. أي مصدر للتشتت يعطي ذلك؟
\end{exercise}

\begin{exercise}[\lvlC]
\label{ex:02:7}
\emph{بحث: قلّاب أم مكثف.} $100$ عصبون تحفظ بتًا مدة $10\,\text{s}$. القلّاب: كل عصبون يطلق $20$ نبضة في الثانية. المكثف: التيسير $u$ يهبط بـ$\tau_F=1\,\text{s}$، ويُقرأ البت عند $u\ge u_{\max}/2$، والإنعاش رشقة من ثلاث نبضات لكل عصبون. كم مرة يكون المكثف أوفر، وماذا يحدث للبت بعد إسكات المجموعة مدة $200\ms$؟
\end{exercise}
